\documentclass[11pt,a4paper]{article}
\usepackage[T1]{fontenc}
\usepackage[utf8]{inputenc}
\usepackage[provide=*,french]{babel}
\usepackage[a4paper,margin=2.3cm]{geometry}
\usepackage{hyperref}
\usepackage{enumitem}
\usepackage{booktabs,longtable,array}
\usepackage{xcolor}
\usepackage{fancyvrb}
\hypersetup{colorlinks=true, linkcolor=blue, urlcolor=blue, pdftitle={Checklist release ZIP --- myco\_apps\_releases\_utilisateurs}}
\providecommand{\tightlist}{\setlength{\itemsep}{0pt}\setlength{\parskip}{0pt}}
\DefineVerbatimEnvironment{Highlighting}{Verbatim}{commandchars=\\\{\}}
\newenvironment{Shaded}{\begin{quote}}{\end{quote}}
\newcommand{\NormalTok}[1]{#1}
\newcommand{\ExtensionTok}[1]{#1}
\setlength{\emergencystretch}{3em}
\title{Checklist release ZIP --- myco\_apps\_releases\_utilisateurs}
\author{Boite Eddy}
\date{\today}
\begin{document}
\maketitle
\section{Checklist release ZIP --- myco\_apps\_releases\_utilisateurs}

\subsection{Objectif}

Préparer une archive autonome, exécutable sans dépendance au dépôt Git.

\subsection{Checklist avant publication}

\begin{enumerate}
\item Vérifier la structure attendue : lanceurs Windows, macOS et Linux, ainsi que le dossier \texttt{applications/}.
\item Vérifier la présence des données d'exemple dans \texttt{applications/data/} et \texttt{exemples/}.
\item Vérifier que \texttt{applications/logs/} est vide (hors \texttt{.gitkeep}).
\item Vérifier que \texttt{applications/results/} est vide (hors \texttt{.gitkeep}).
\item Supprimer les artefacts système (ex. \texttt{.DS\_Store}).
\item Tester un lancement rapide sur les scripts de lancement des OS cibles.
\item Vérifier qu'un run ICR et un run CHEGD produisent bien des sorties dans le dossier de résultats.
\item Vérifier qu'un log d'exécution est bien créé dans \texttt{applications/logs/}.
\item Vérifier que la documentation utilisateur est cohérente avec les noms de fichiers réels.
\item Générer le ZIP uniquement à partir du dossier de release utilisateur.
\end{enumerate}

\subsection{Contrôle final}

La release est valide si un utilisateur peut :
\begin{itemize}
\item dézipper,
\item lancer l'application,
\item exécuter les deux analyses avec les exemples,
\item récupérer les résultats et les logs localement.
\end{itemize}
\end{document}